\chapter{Batch Clearing Algorithm}
\label{app:algorithm}

The clearing routine implements \eqref{eq:demandsupply} and
\eqref{eq:clearing} by a search over the admissible price grid, followed
by the tie-break cascade and the rationing step of
\Cref{sec:rationing}.

\begin{lstlisting}[language=,caption={Batch clearing, single asset
pair.},label={lst:clearing}]
INPUT : buys B = {(pi_i, q_i)}, sells S = {(pi_j, q_j)},
        tick grid P, tie-break cascade C, rationing rule R,
        previous clearing price p_prev
OUTPUT: clearing price p*, fills F, residual E

1  candidates <- distinct limit prices in B and S, snapped to P
2  for each p in candidates:
3      D(p) <- sum of q_i over buys with pi_i >= p
4      S(p) <- sum of q_j over sells with pi_j <= p
5      V(p) <- min(D(p), S(p))
6  Vmax  <- max over p of V(p)
7  M     <- { p : V(p) = Vmax }              // maximising set
8  if Vmax = 0: return (no clearing, all orders roll over)
9  M     <- argmin over p in M of |D(p) - S(p)|        // cascade 1
10 if |M| > 1: M <- argmin over p in M of |p - p_prev|  // cascade 2
11 if |M| > 1: M <- argmin over p in M of p             // cascade 3
12 p*    <- unique element of M
13 fill all buys with pi_i > p* completely            // inframarginal
14 fill all sells with pi_j < p* completely
15 marginal <- orders with pi = p* on the long side
16 allocate Vmax - (filled inframarginal) among marginal using R
17 E     <- unmatched quantity on the long side
18 apply time-in-force to E: roll over, or cancel
19 return (p*, F, E)
\end{lstlisting}

Steps 1--7 are the line search; the unimodality of $V$ noted in
\Cref{sec:clearing} guarantees that the maximising set is an interval on
the grid. Steps 9--11 are the tie-break cascade: minimise the residual
imbalance, then minimise distance from the auction's own previous
clearing price $p_{\mathrm{prev}}$ (not an external reference — the FBA
has no oracle input, see \Cref{app:config}), then take the lowest price
as a final deterministic tie-break. Step 16 is the only place where the
rationing rule enters, which is why the rule can be swapped without
touching the price computation. Steps 13--16 are written as two phases
for clarity; the implementation instead walks each side once in a single
priority order, inframarginal orders sorted ahead of marginal ones, so
steps 13--14's ``completely'' holds in the typical case but not
unconditionally: when the maximising interval spans more than one price,
a strictly better-priced order on the \emph{longer} side can itself be
left only partly filled (\Cref{sec:testfba}).
